\documentclass[./main.tex]{subfiles}

\begin{document}

\begin{center}
    {\Large\bfseries Utility Billing for Industrial Energy Assessment}
\end{center}

\section*{Model 1: Two Facilities, Same Annual Energy, Different Bills}

The table below shows annual electricity billing data for two manufacturing
facilities. Both consumed the same total kWh over the year, but their bills
look very different. Study the data carefully before answering the questions.

\begin{longtable}[]{@{}
  >{\raggedright\arraybackslash}p{(\columnwidth - 4\tabcolsep) * \real{0.4060}}
  >{\raggedright\arraybackslash}p{(\columnwidth - 4\tabcolsep) * \real{0.2970}}
  >{\raggedright\arraybackslash}p{(\columnwidth - 4\tabcolsep) * \real{0.2970}}@{}}
\toprule\noalign{}
Billing item & Acme Stamping & Riverside Plastics \\
\midrule\noalign{}
\endhead
\bottomrule\noalign{}
\endlastfoot
\multicolumn{3}{@{}l@{}}{\textbf{Annual consumption}} \\
Total kWh consumed    & 1,200,000 kWh            & 1,200,000 kWh \\
Operating hours       & 8,760 hr/yr (continuous)  & 2,000 hr/yr (day shift only) \\
\multicolumn{3}{@{}l@{}}{\textbf{Peak demand}} \\
Peak 15-min demand    & 145 kW                    & 720 kW \\
\multicolumn{3}{@{}l@{}}{\textbf{Annual charges}} \\
Annual consumption charges    & \$84,000      & \$84,000 \\
Annual demand charges         & \$20,880      & \$103,680 \\
Total annual electricity cost & \$104,880     & \$187,680 \\
Blended rate (total cost / total kWh) & \$0.0874/kWh & \$0.1564/kWh \\
\end{longtable}

\begin{enumerate}
  \needspace{4\baselineskip}
  \item Looking only at the total kWh consumed row, what do you notice about
        the two facilities? Why is this the case?
        \vspace{1.25in}

  \needspace{4\baselineskip}
  \item Compare the annual demand charges. Which facility pays more, and by
        roughly how much?
        \vspace{1.25in}

  \needspace{4\baselineskip}
  \item Explain in your own words why Riverside's demand charges are so much
        higher than Acme's, even though their annual kWh consumption is
        identical.
        \vspace{1.25in}

  \needspace{4\baselineskip}
  \item Acme runs 24 hours a day, 365 days a year. Riverside runs only a day
        shift ($\sim$2,000 hours per year). Both used the same total kWh.
        Calculate the average power level of each facility during operating
        hours using the formula below, then record your results.

        \[
          \text{Average Power (kW)} = \frac{\text{Total kWh Consumed}}{\text{Operating Hours (hr/yr)}}
        \]
        \vspace{1.25in}

  \needspace{4\baselineskip}
  \item Load factor is defined as the quotient of average load and peak load.
        Calculate the load factor for each facility.

        \[
          \text{Load Factor} = \frac{\text{Average Load (kW)}}{\text{Peak Load (kW)}}
        \]
        \vspace{1.25in}

  \needspace{4\baselineskip}
  \item If two facilities have the same annual consumption, will the one with
        higher or lower load factor pay more in demand charges?
        \vspace{1.25in}

  \needspace{4\baselineskip}
  \item Demand charges are typically assessed on the peak 15-minute interval
        in the billing period. Why might Riverside Plastics' peak demand be
        720 kW while their average demand is lower?
        \vspace{1.25in}

  \needspace{4\baselineskip}
  \item Look at the blended rate row. The two values are different even though
        both facilities pay the same tariff rates. Explain why the blended
        rates differ.
        \vspace{1.25in}
\end{enumerate}

\clearpage
\section*{Model 2: The Blended Rate Trap}

The blended rate seems convenient: divide the total bill by total kWh to get a
single number. The table below shows what happens when two efficiency measures
are evaluated using the blended rate versus the correct component rates.

\medskip
\noindent\textbf{Facility rate structure for this model:}
\begin{itemize}
      \item Consumption rate: \$0.07/kWh
      \item Demand rate: \$12/kW/month
      \item Blended rate: \$0.115/kWh
\end{itemize}


\begin{longtable}[]{@{}
  >{\raggedright\arraybackslash}p{(\columnwidth - 4\tabcolsep) * \real{0.2991}}
  >{\raggedright\arraybackslash}p{(\columnwidth - 4\tabcolsep) * \real{0.3504}}
  >{\raggedright\arraybackslash}p{(\columnwidth - 4\tabcolsep) * \real{0.3504}}@{}}
\toprule\noalign{}
\textbf{Item}
& \textbf{Measure A: Lighting Controls (saves 100,000 kWh/yr; no demand reduction)}
& \textbf{Measure B: Load scheduling (saves 0 kWh/yr; reduces peak demand by 50 kW)} \\
\midrule\noalign{}
\endhead
\bottomrule\noalign{}
\endlastfoot
\multicolumn{3}{@{}l@{}}{\textbf{Correct calculation (component rates)}} \\
Consumption savings
  & 100,000 kWh $\times$ \$0.07/kWh $=$ \$7,000/yr
  & 0 kWh $\times$ \$0.07/kWh $=$ \$0/yr \\
Demand savings
  & \$0 (no peak reduction)
  & 50 kW $\times$ \$12/kW $\times$ 12 mo $=$ \$7,200/yr \\
\textbf{Correct total savings} & \textbf{\$7,000/yr} & \textbf{\$7,200/yr} \\
\multicolumn{3}{@{}l@{}}{\textbf{Blended rate calculation}} \\
Blended rate estimate
  & 100,000 kWh $\times$ \$0.115/kWh $=$ \$11,500/yr
  & 0 kWh $\times$ \$0.115/kWh $=$ \$0/yr \\
\textbf{Blended rate error}
  & \textbf{Overstated by \$4,500 (+64\%)}
  & \textbf{Understated by \$7,200 ($-$100\%)} \\
\end{longtable}

\begin{enumerate}
  \setcounter{enumi}{8}
  \needspace{4\baselineskip}
  \item For Measure A (lighting upgrade), the blended rate overstated savings
        by 64\%. Explain mechanically why this happens. Where does the extra
        \$4,500 come from?
        \vspace{1.25in}

  \needspace{4\baselineskip}
  \item For Measure B (load scheduling), the blended rate returned \$0 even
        though the correct calculation yields \$7,200/yr. Why does the blended
        rate fail completely for this measure?
        \vspace{1.25in}

  \needspace{4\baselineskip}
  \item Based on the two failure modes in Model 2, complete the following
        statement:

        \vspace{0.25in}
      The blended rate embeds \underline{\hspace{4cm}} charges into a
      per-kWh number, which causes it to overstate savings for measures that
      only reduce \underline{\hspace{4cm}} and completely miss savings for
      measures that only reduce \underline{\hspace{4cm}}.
      \vspace{0.25in}


  \needspace{4\baselineskip}
  \item A facilities employee at a major defense contractor says: ``I always
        use the blended rate because it's simpler and the errors cancel out in
        the end.'' Evaluate this claim. Under what facility profile would the
        errors actually cancel?
        \vspace{1.25in}

  \needspace{4\baselineskip}
  \item At ITAC, we never use blended rates. Explain the rationale for that
        policy based on what you observed in Model 2.
        \vspace{1.25in}
\end{enumerate}

\clearpage
\section*{Model 3: Demand Ratchets and the True Cost of a Peak}

Many commercial and industrial tariffs include a ratchet clause. The diagram
below shows how a ratchet operates over a 12-month window for a facility that
had an unusually high demand spike in January. Their demand rate is \$12/kW.

\medskip
\noindent\textbf{Tariff terms:} Billed demand $=$ max(actual monthly demand,
$0.80 \times$ highest demand in preceding 12 months)
\medskip

\begin{figure}[ht]
  \centering
  \fbox{\parbox{0.85\textwidth}{\centering\vspace{3cm}
    \textit{[Figure 1: Demand ratchet bar chart --- insert image here]}
  \vspace{3cm}}}
  \caption{Monthly billed demand vs.\ actual demand under an 80\% ratchet
    clause following a January demand spike.}
\end{figure}

\begin{enumerate}
  \setcounter{enumi}{13}
  \needspace{4\baselineskip}
  \item In which months does the ratchet clause result in billed demand
        exceeding actual demand?
        \vspace{1.25in}

  \needspace{4\baselineskip}
  \item For March (actual demand = 290 kW), calculate the demand charge both
        without and with the ratchet applied. Show your work. What is the
        additional monthly cost imposed by the ratchet?
        \vspace{1.25in}

  \needspace{4\baselineskip}
  \item The January spike was caused by an unusually cold week that required
        all electric heating systems to run simultaneously. Suggest two
        operational strategies that could have prevented the spike from
        reaching 500 kW.
        \vspace{1.25in}

  \needspace{4\baselineskip}
  \item An equipment upgrade would reduce the facility's peak demand by 70 kW
        in every month. The measure costs \$10,000. Using \$12/kW/month:
        \begin{enumerate}
          \item Calculate the annual demand savings from the 70 kW reduction
                alone (ignoring the ratchet clause).
                \vspace{1.25in}
          \item Now consider that the 80\% ratchet is still in effect and the
                500 kW historical peak has not yet rolled off the 12-month
                window. Does the demand reduction produce any bill savings
                immediately? Explain.
                \vspace{1.25in}
          \item What does this tell you about the timing of demand-reduction
                savings at facilities with ratchet clauses?
                \vspace{1.25in}
        \end{enumerate}
\end{enumerate}

\clearpage
\section*{Model 4: Power Factor and Its Effect on the Bill}

Power factor (PF) describes how efficiently a facility converts the electricity
delivered by the utility into useful work. The triangle below shows the
relationship between real power, reactive power, and apparent power.

\begin{figure}[ht]
  \centering
  \fbox{\parbox{0.85\textwidth}{\centering\vspace{3cm}
    \textit{[Figure 2: Power factor triangle diagram --- insert image here]}
  \vspace{3cm}}}
  \caption{Power triangle illustrating the relationship between real power
    (kW), reactive power (kVAR), and apparent power (kVA).}
\end{figure}

The table below shows the power factor of a facility before and after they
install capacitor banks to increase their power factor. Capacitor banks help
with this by supplying leading reactive power, which cancels out lagging
reactive power.

\begin{longtable}[]{@{}
  >{\raggedright\arraybackslash}p{(\columnwidth - 4\tabcolsep) * \real{0.3745}}
  >{\raggedright\arraybackslash}p{(\columnwidth - 4\tabcolsep) * \real{0.2858}}
  >{\raggedright\arraybackslash}p{(\columnwidth - 4\tabcolsep) * \real{0.3397}}@{}}
\toprule\noalign{}
\textbf{Parameter} & \textbf{Before correction} & \textbf{After correction} \\
\midrule\noalign{}
\endhead
\bottomrule\noalign{}
\endlastfoot
\multicolumn{3}{@{}l@{}}{\textbf{Electrical quantities}} \\
Real power (kW)       & 200 kW             & 200 kW \\
Apparent power (kVA)  & 267 kVA            & 211 kVA \\
Reactive power (kVAR) & 177 kVAR (lagging) & 66 kVAR (lagging) \\
kVAR correction needed for 0.95 PF & --- & 177 kVAR $-$ 66 kVAR $=$ 111 kVAR correction \\
\end{longtable}

\begin{enumerate}
  \setcounter{enumi}{17}
  \needspace{4\baselineskip}
  \item Calculate the power factor before and after correction. Show your
        calculations.
        \vspace{1.25in}

  \needspace{4\baselineskip}
  \item The real power consumed by the facility did not change before or after
        power factor correction. But the apparent power dropped significantly.
        Explain physically what the capacitor bank is doing to make this
        possible. (\emph{Hint: draw the power factor triangle before and after
        correction --- what happens to the triangle?})
        \vspace{1.25in}

  \needspace{4\baselineskip}
  \item Utilities must generate, transmit, and distribute apparent power
        (kVA), but they can only bill for real power (kW) on some tariffs.
        Why does this create a financial incentive for utilities to penalize
        low power factor?
        \vspace{1.25in}

  \needspace{4\baselineskip}
  \item For each of the rate tariffs in the table below, calculate the demand
        savings from the scenario illustrated in Model 4 by calculating the
        billed demand before/after the correction, then finding the difference.
        The facility is billed at \$12 per kW or kVA.

        \begin{longtable}[]{@{}
          >{\raggedright\arraybackslash}p{(\columnwidth - 4\tabcolsep) * \real{0.3333}}
          >{\raggedright\arraybackslash}p{(\columnwidth - 4\tabcolsep) * \real{0.3334}}
          >{\raggedright\arraybackslash}p{(\columnwidth - 4\tabcolsep) * \real{0.3334}}@{}}
        \toprule\noalign{}
        Rate Tariff 1 & Rate Tariff 2 & Rate Tariff 3 \\
        \midrule\noalign{}
        \endhead
        \bottomrule\noalign{}
        \endlastfoot
        Demand is billed as the greater of real power or 90\% of apparent power
        & If power factor falls below 90\%, the company's demand is billed by kVA
          instead of kW.
        & Demand is billed by kVA, regardless of power factor. \\
        \end{longtable}
        \vspace{1.25in}

  \needspace{4\baselineskip}
  \item Based on the power factor post-correction, which rate tariff do you
        think the company was under? Why?
        \vspace{1.25in}
\end{enumerate}

\end{document}
